% !TEX root = ../frode-bibelen.tex
% ----------------------------------------------------------------------------
%  Frode-bibelen -- 11 Oppgave 1: Hva blir utskriften?
%  Inkluderes fra frode-bibelen.tex (kompiler alltid hovedfilen).
% ----------------------------------------------------------------------------

% ===================================================== 11 OPPGAVE 1 =======
\kapittel{Oppgave 1: «Hva blir utskriften?»}

\begin{regel}[{Metode -- gjør det mekanisk, ikke i hodet}]
\begin{enumerate}
  \item \textbf{Skriv teksten(e) ut med indekser under} (bruk sporingsarket, vedlegg B). Da blir \cc{txt[i]}, \cc{txt+13} og \cc{*(t+2)} rene oppslag.
  \item \textbf{Lag en variabeltabell} og oppdater den for hver setning -- spesielt pekere (skriv hvilken \emph{indeks} pekeren står på).
  \item \textbf{Én \cc{printf} om gangen}: noter mellomrom, \cc{\%5i}-bredder og \cc{\bs n} nøyaktig.
  \item \textbf{Tell linjene}: oppgaven sier hvor mange linjer det blir. Stemmer ikke antallet, har du en feil.
  \item \textbf{Funksjonskall}: får funksjonen en \emph{kopi} (endringen forsvinner) eller en \emph{peker} (originalen endres)?
  \item \textbf{Løkkegrenser}: \cc{<} vs.\ \cc{<=}, \cc{i += 8}, start/slutt -- skriv opp alle verdiene \cc{i} tar.
\end{enumerate}
\end{regel}

\subsection{Oppslag: hva betyr uttrykket?}
\begin{center}
\begin{longtable}{@{}L{3.9cm}L{11.7cm}@{}}
\toprule
\textbf{Uttrykk} & \textbf{Betydning} \\
\midrule
\endhead
\cc{t = txt;} & \cc{t} peker på indeks 0. \cc{t = txt+5;} peker på indeks 5. \\
\cc{*t}, \cc{t[0]} & Tegnet der \cc{t} peker. \\
\cc{*(t+2)}, \cc{t[2]} & Tegnet \textbf{2 plasser etter} \cc{t} -- \cc{t} flyttes \emph{ikke}. \\
\cc{*t+1} & \cc{(*t)+1}: tegnet \emph{pluss én} i ASCII (\cc{'a'+1} = \cc{'b'}). \\
\cc{*t++} & Bruk \cc{*t}, flytt \cc{t} \textbf{etterpå} (også når sammenligningen er usann!). \\
\cc{*++t} & Flytt \cc{t} først, bruk så det nye tegnet. \\
\cc{++t;} / \cc{t++;} & Alene på en linje: begge flytter \cc{t} én plass. \\
\cc{while (*t)} & Til \cc{t} står på \cc{'\bs 0'} (slutten). \\
\cc{printf("\%s", t)} & Skriver fra \cc{t} \textbf{til slutten} av teksten. \\
\cc{printf("\%c", t)} vs.\ \cc{*t} & \cc{\%c} trenger et \emph{tegn}: \cc{*t} eller \cc{t[i]}. \\
\cc{strchr(s, c)} & Peker til første \cc{c} (eller \cc{NULL}). \cc{strchr(++t, c)} / \cc{strchr(t+1, c)} finner \textbf{neste}. \\
\cc{strstr(s, "al")} & Peker til første \cc{"al"} (eller \cc{NULL}). Etterpå ofte \cc{t += 2} for å gå forbi. \\
\cc{!strncmp(a, b, n)} & \textbf{Sann (1)} når de \cc{n} første tegnene er \emph{like}. \cc{!strcmp} tilsvarende. \\
\cc{strcmp(a, b) < 0} & \cc{a} kommer før \cc{b} i ASCII (store bokstaver før små!). \\
\cc{strcpy(t, s)} / \cc{strcat(t, s)} & Overskriv / skjøt på. \\
\cc{printf("\%i", a > b)} & Sannhetsverdier skrives som \cc{1} eller \cc{0}. Det samme for \cc{bool}. \\
\cc{(char)(*c + 2)} & To plasser lenger ut i ASCII-tabellen (\cc{'M'} $\rightarrow$ \cc{'O'}). Se vedlegg A. \\
\cc{c - 'a'}, \cc{c - '0'} & Plass i alfabetet / sifferverdi. \\
\cc{7/2}, \cc{7\%2}, \cc{7/2.0} & \cc{3}, \cc{1}, \cc{3.5}. \ \ \cc{7/2*2} = \cc{6}. \\
\cc{n *= 2} i løkke & 1, 2, 4, 8, 16 \dots\ -- skriv opp verdien \emph{før} og \emph{etter} hver \cc{printf}. \\
\cc{a[n++] = x;} & Lagrer på plass \cc{n}, øker \cc{n} etterpå. \\
\cc{char txt[][30]} & \cc{txt[i]} er tekst nr.\ \cc{i}, \cc{txt[i][j]} er tegn nr.\ \cc{j} i den. \\
\cc{void f(struct S s)} & Kopi: endringer (også \cc{strcat} på \cc{s.navn[]}) forsvinner. \\
\cc{void f(struct S* s)} & Original: endringer blir stående. \\
\cc{a = b;} (structer) & Kopierer alt innhold, også \cc{char}-arrayer inni. \\
\cc{if (s)}, \cc{if (p)} & Sann når \cc{s} ikke er 0 / \cc{p} ikke er \cc{NULL}. \\
\cc{\&\&}, \cc{||} & Kortslutning: høyre side kjøres kanskje \emph{ikke} (viktig hvis den har \cc{++}). \\
\bottomrule
\end{longtable}
\end{center}

\subsection{Gjennomregnet eksempel A (tekst og pekere)}
\begin{kode}
char txt[] = "Gjovik-Hamar-Lillehammer";

int main()  {
  char* t = txt;
  int n = 0;

  for (int i = 0;  i < 24;  i += 5)  printf("%c", txt[i]);
  printf("\n");

  while (*t)  {  if (*t++ == 'a')  n++;  }
  printf("%i %s\n", n, txt+13);

  t = strchr(txt, '-');
  printf("%c%c %c\n", *(t+1), *(t-1), (char)(*t + 1));

  printf("%i %i %.1f\n", !strncmp(txt+7, "Ham", 3), 7 / 2 * 2, 7 / 2.0);
  return 0;
}
\end{kode}

\textbf{Trinn 1 -- indekser:}\par
{\ttfamily\footnotesize\setlength{\tabcolsep}{2.2pt}
\begin{tabular}{|*{24}{c|}}
\hline
G&j&o&v&i&k&-&H&a&m&a&r&-&L&i&l&l&e&h&a&m&m&e&r\\ \hline
0&1&2&3&4&5&6&7&8&9&10&11&12&13&14&15&16&17&18&19&20&21&22&23\\ \hline
\end{tabular}}

\textbf{Trinn 2 -- linje for linje:}
\begin{enumerate}
  \item \cc{i} = 0, 5, 10, 15, 20 $\Rightarrow$ \cc{G k a l m} $\Rightarrow$ \cc{Gkalm}
  \item \cc{'a'} står på 8, 10, 19 $\Rightarrow$ \cc{n = 3}. \ \cc{txt+13} er \cc{"Lillehammer"} $\Rightarrow$ \cc{3 Lillehammer}
  \item \cc{t} på indeks 6 (\cc{'-'}). \cc{*(t+1)} = \cc{'H'}, \cc{*(t-1)} = \cc{'k'}, \cc{'-'+1} = \cc{'.'} (ASCII 45+1 = 46) $\Rightarrow$ \cc{Hk .}
  \item Indeks 7--9 er \cc{"Ham"} $\Rightarrow$ 1. \ \cc{7/2*2} = \cc{3*2} = 6. \ \cc{7/2.0} = 3.5 $\Rightarrow$ \cc{1 6 3.5}
\end{enumerate}

\subsection{Gjennomregnet eksempel B (struct og parametre)}
\begin{kode}
struct Person  {  char navn[20];  int alder;  };

void funk1(struct Person p)    {  p.alder++;   strcat(p.navn, "!");  }
void funk2(struct Person* p)   {  p->alder++;  strcat(p->navn, "?");  }
bool funk3(const struct Person* p, int n)  {  return (strlen(p->navn) > n);  }

int main()  {
  struct Person a = { "Ola", 20 },  b = { "Kari", 30 };

  funk1(a);   funk2(&b);
  printf("%s %i - %s %i\n", a.navn, a.alder, b.navn, b.alder);

  a = b;   funk2(&a);
  printf("%s %i - %s %i\n", a.navn, a.alder, b.navn, b.alder);

  printf("%i %i\n", funk3(&a, 5), funk3(&b, 5));
  return 0;
}
\end{kode}
\begin{enumerate}
  \item \cc{funk1} får en \textbf{kopi} av \cc{a} $\Rightarrow$ \cc{a} uendret. \cc{funk2} får \textbf{adressen} til \cc{b} $\Rightarrow$ \cc{b} blir \cc{"Kari?"}, 31.\\
        Utskrift: \cc{Ola 20 - Kari? 31}
  \item \cc{a = b} kopierer alt (\cc{"Kari?"}, 31). \cc{funk2(\&a)} endrer bare \cc{a}.\\
        Utskrift: \cc{Kari?? 32 - Kari? 31}
  \item \cc{strlen("Kari??")} = 6 > 5 $\Rightarrow$ 1. \ \cc{strlen("Kari?")} = 5 > 5 er usant $\Rightarrow$ 0.\\
        Utskrift: \cc{1 0}
\end{enumerate}

\begin{tips}[{Øv på ekte oppgaver}]
Alle oppgave 1-ene fra H19 til S26 ligger i \href{https://frh.folk.ntnu.no/grprog/eksamen/}{eksamensarkivet}
(\cc{ex\_xxx\_1a.c} og \cc{ex\_xxx\_1b.c}, med fasit nederst i filen).
Løs dem med penn og sporingsark \emph{før} du ser på fasiten -- og kjør dem gjerne i Code::Blocks etterpå.
\end{tips}
